\subsection{Phân tích thị trường và đối thủ cạnh tranh}
\label{sec:market-research}

\begin{guide}
\guidepart{Nội dung}
\begin{compactitem}
    \item Đối thủ trực tiếp (ứng dụng giải trí tổng hợp) và gián tiếp (ứng dụng chuyên biệt cho từng phân hệ: video, nhạc, tin tức, game).
    \item Khảo sát định tính (phỏng vấn) và định lượng (bảng khảo sát): mục tiêu, số người tham gia, câu hỏi chính, kết quả tổng hợp.
    \item Ma trận SWOT (Strengths, Weaknesses, Opportunities, Threats) của ứng dụng.
\end{compactitem}
\guidepart{Câu hỏi cần trả lời}
\begin{compactitem}
    \item Đối thủ nào đang được dùng nhiều nhất? Họ làm tốt điều gì và còn thiếu điều gì?
    \item Khảo sát được thực hiện với bao nhiêu người, theo hình thức nào? Kết quả nào ảnh hưởng trực tiếp đến thiết kế?
    \item Từ SWOT, điểm mạnh nào cần khai thác và rủi ro nào cần có phương án giảm thiểu?
\end{compactitem}
\guidepart{Hình/bảng nên có}
Bảng so sánh tính năng với đối thủ, biểu đồ kết quả khảo sát, ma trận SWOT $2 \times 2$.
\end{guide}

% Ghi chú phân tích đối thủ do nhóm tổng hợp (dành cho mục 1.3, chưa biên tập thành nội dung):
% Nhóm đối thủ gián tiếp (ứng dụng chuyên biệt - đang chiếm thị phần lớn nhất):
% - Video Streaming (YouTube / TikTok): làm tốt kho nội dung đồ sộ, thuật toán gợi ý chuẩn xác,
%   luồng stream mượt mà; còn thiếu phát nhạc nền chuyên sâu nếu không có gói Premium, không có
%   không gian đọc tin tức văn bản và minigame tương tác trên cùng ứng dụng.
% - Music Listening (Spotify / Zing MP3): làm tốt chất lượng âm thanh, cá nhân hoá playlist,
%   chạy nền mượt mà; còn thiếu đọc tin tức, xem video dài đa dạng và gamification/minigame.
% - News Reading (VnExpress / Báo Mới): làm tốt tốc độ cập nhật, giao diện đọc báo tối ưu;
%   còn thiếu trải nghiệm nghe nhạc đệm khi đọc tin, nội dung chủ yếu văn bản và ảnh tĩnh.
% - Minigames (Shopee Games / game casual lẻ): làm tốt giải trí nhanh, tích điểm tạo động lực;
%   còn thiếu gắn với hệ sinh thái truyền thông, tồn tại rời rạc hoặc phụ thuộc ứng dụng TMĐT.
